\documentclass{article}
\usepackage[T1]{fontenc}
\usepackage{lmodern}
\usepackage{graphicx}
\usepackage{amsmath,amssymb}
\usepackage[a4paper,margin=2cm]{geometry}
\usepackage[absolute,overlay]{textpos}
\setlength{\TPHorizModule}{1mm}
\setlength{\TPVertModule}{1mm}
\usepackage[indonesian]{babel}
\usepackage{hyperref}
\setlength{\emergencystretch}{2em}
% unit: Halaman 30

\begin{document}

% === Halaman 30 (llm) ===

\section*{Dekripsi}

\begin{itemize}
    \item Dekripsi terhadap cipherteks merupakan kebalikan dari proses enkripsi.
    \item DES menggunakan algoritma yang sama untuk proses enkripsi dan dekripsi.
    \item Pada proses dekripsi urutan kunci yang digunakan adalah $K_{16}, K_{15}, \dots, K_1$.
    \item Untuk tiap putaran $16, 15, \dots, 1$, luaran pada setiap putaran \textit{deciphering} adalah
    \[
    \begin{aligned}
    R_{i-1} &= L_i \\
    L_{i-1} &= R_i \oplus f(R_{i-1}, K_i) = R_i \oplus f(L_i, K_i)
    \end{aligned}
    \]
\end{itemize}

\end{document}
